\documentclass[10pt, twoside]{article}

%% Стилевой файл, обеспечивающий класс LaTeX2e для журнала
%% кафедры математической экономики и эконометрики НИУ ВШЭ
%% «Эпсилон» («Epsilon»)
%%
%% Данный стилевой файл распространяется свободно и без ограничений
%% при условии, что в него не будет вноситься изменений.
%% Если изменения ПРИСУТСТВУЮТ, то вы должны переименовать этот файл.
%%
%%  This work may be distributed and/or modified under the
%%  conditions of the LaTeX Project Public License, either version 1.3
%%  of this license or (at your option) any later version.
%%  The latest version of this license is in
%%    http://www.latex-project.org/lppl.txt
%%  and version 1.3 or later is part of all distributions of LaTeX
%%  version 2005/12/01 or later.
%%
%%  This work has the LPPL maintenance status `maintained'.
%%
%%  The Current Maintainer of this work is A. V. Kostyrka.
%%
%%  This work consists of the file epsilonj.sty
%%
%% Bug fixes and changes:
%% TeX programming: Andreï V. Kostyrka, HSE, Russia, avkostyrka@edu.hse.ru
%% Requires Latex2e, ver.2000.06

\NeedsTeXFormat{LaTeX2e}
\ProvidesPackage{epsilonj}[2014/04/21 by Andre\"{\i} Kostyrka (v. 0.1)]

\RequirePackage[utf8]{inputenc}
\RequirePackage[X2,T2A]{fontenc}
\RequirePackage[english, russian, KeepShorthandsActive]{babel}
\RequirePackage[cmex10]{amsmath}
\RequirePackage{amsfonts,amssymb,amsthm}
\RequirePackage{xspace}
\RequirePackage{xcolor}
\RequirePackage{booktabs}
\RequirePackage{indentfirst}
\RequirePackage{microtype}
\RequirePackage{graphicx}
\RequirePackage{subfigure}
\RequirePackage[margin=0pt,font=small,labelfont=bf,labelsep=period]{caption}
\RequirePackage{titling}
\RequirePackage{ifthen} % For checking ifx with empty commands

\RequirePackage[textwidth=32pc, textheight=48pc, paperwidth=15cm, paperheight=24cm]{geometry}

\RequirePackage{enumitem}
\setlist{itemsep=0ex plus.1ex, parsep=0ex plus.1ex, topsep=.7ex plus.2ex minus.15ex, partopsep=.7ex plus.2ex minus.15ex}
\AddEnumerateCounter{\Asbuk}{\@Asbuk}{\CYRM}
\AddEnumerateCounter{\asbuk}{\@asbuk}{\cyrm}
\renewcommand{\labelenumii}{\theenumii)}
\renewcommand{\theenumii}{\asbuk{enumii}}


\RequirePackage{csquotes}
\DeclareQuoteStyle{russian}
{\guillemotleft}
{\guillemotright}
{\quotedblbase}
{\textquotedblright}

\RequirePackage[
  backend=biber,
  natbib=true,
  bibstyle=gost-authoryear,
  citestyle=gost-authoryear,
  movenames=false,
  sortlocale=ru_RU,
  url=false,
  doi=false,
  isbn=false,
  eprint=false,
  language=auto
]{biblatex}
% \addbibresource{epsilon.bib}
%\usepackage{paratype}

\RequirePackage[colorlinks,citecolor=blue,urlcolor=blue]{hyperref}

%%%%%%%%%%%%%%%%%%%%%%%%%%%%%%%%%%%%%% Modifications and "add-on" for article.cls starts:

%%%%%%%%%%%%%%%%%%%%%%%%%%%%%%%%%%%%%% Initiate some info:

\usepackage{fancyhdr}
\setlength{\headheight}{15.2pt}

\pagestyle{fancy}
%\renewcommand{\sectionmark}[1]{ \markright{#1}{} }

\newcommand{\fullAUTH}{}
\newcommand{\runAUTH}{}
\newcommand{\fullTITL}{}
\newcommand{\runTITL}{}

\newenvironment{keyword}
{\textbf{Ключевые слова:}}
{\unskip.}

\fancyhf{}
\fancyhead[LE, RO]{\thepage}
\fancyhead[CE]{\runAUTH}
\fancyhead[CO]{\runTITL}

\newcommand{\AUTHOR}[2]{%
	\ifx&#1&%
	\renewcommand{\fullAUTH}{#1} \author{}%
    \else%
	    \ifx&#2&%
		\renewcommand{\fullAUTH}{#1} \author{#1}%
		\fi%
	\renewcommand{\fullAUTH}{#1} \author{#1\thanks{#2}}%
    \fi%
}
\newcommand{\SHORTAUTHOR}[1]{\renewcommand{\runAUTH}{#1}}
\newcommand{\TITLE}[1]{\renewcommand{\fullTITL}{#1} \title{#1}}
\newcommand{\SHORTTITLE}[1]{\renewcommand{\runTITL}{#1}}

\fancypagestyle{firststyle}
{
   \fancyhf{}
   \fancyhead[C]{\href{http://economics.hse.ru/depe/econometrics/}{Опубликовано в журнале «Эпсилон»}}
}

\newcommand{\DoFirstPageTechnicalStuff}{%
	\newread\inputstream
	\openin\inputstream=0pagecount.txt
	\read\inputstream to \PageCount
	\closein\inputstream
	\setcounter{page}{\PageCount}%
	\thispagestyle{firststyle}%
	\maketitle%
	}
	
\newcommand{\DoFirstPageSimple}{%
	\newread\inputstream
	\openin\inputstream=0pagecount.txt
	\read\inputstream to \PageCount
	\closein\inputstream
	\setcounter{page}{\PageCount}%
	\thispagestyle{firststyle}%
}

% % % % % % % % % % % % % % % % % % % % % % % User command (AVK)

\def\ellipsiskern{.1em\relax}
\newcommand{\ldotst}{.\kern\ellipsiskern.\kern\ellipsiskern.} % для . . .
\newcommand{\ldotse}{!\kern\ellipsiskern.\kern\ellipsiskern.} % для ! . .
\newcommand{\ldotsq}{?\kern\ellipsiskern\kern-.11em.\kern\ellipsiskern.} % для ? . .
\newcommand{\fnnsp}{\kern-0.4em\relax}

\def\tagform@#1{\maketag@@@{{\footnotesize\textsf{(#1)}}\@@italiccorr}}

\def\hm#1{#1\nobreak\discretionary{}{\hbox{$#1$}}{}}

\newcommand{\E}{\mathbb{E}}                 % Мат. ожидание
\newcommand{\e}{\varepsilon}                % эпсилон
\newcommand{\p}{\partial}                   % частная производная
\newcommand{\iu}{{\mkern1mu\mathrm{i}\mkern1mu}} % Мнимая единица
\newcommand{\ee}{{\mkern1mu\mathrm{e}\mkern1mu}} % Экспонента
\newcommand{\PP}{\mathbb{P}}                % Вероятность
\newcommand{\bs}[1]{\boldsymbol{#1}}        % ЖИРНЫЙ последующий символ имени матрицы или векторы
\renewcommand{\to}{\rightarrow}                   % Правильная стрелка вправо («стремится»)
\newcommand{\const}{\ensuremath{\text{const}}}        % const прямым начертанием
\newcommand{\CONST}{\mathbb{C}}
\newcommand{\assim}{\mathrel{\stackrel{\text{as}}\sim}}
\newcommand{\iid}{\text{i.\kern1pt i.\kern1pt d.}}
\renewcommand{\iff}{\hskip.5em plus.1em minus.1em \Leftrightarrow\hskip.5em plus.1em minus.1em\relax}
\newcommand{\protivor}{\text{\textreferencemark}}
\providecommand{\hence}{\hskip.5em plus.1em\Rightarrow\hskip.5em plus.1em}
\newcommand{\ind}{\mathbf{1}}
\newcommand{\ddd}{\mathrm{d}}
\renewcommand{\emptyset}{\varnothing}
\newcommand{\calN}{\mathcal{N}}
\newcommand{\calL}{\mathcal{L}}
\newcommand{\simhypo}{\mathrel{\stackrel{\hypo_0}{\sim}}}
\newcommand{\ofe}[1]{\bigl( \{#1\} \bigr)}
\DeclareMathOperator{\Corr}{Corr}       % Оператор корреляции
\DeclareMathOperator{\corr}{Corr}
\DeclareMathOperator{\Cov}{Cov}         % Оператор ковариации
\DeclareMathOperator{\cov}{Cov}
\DeclareMathOperator{\Var}{Var}         % Оператор дисперсии
\DeclareMathOperator{\var}{Var}
\DeclareMathOperator{\rang}{rank}       % Оператор ранга
\DeclareMathOperator{\rank}{rank}
\DeclareMathOperator*{\plim}{plim}      % Оператор предела по вероятности
\DeclareMathOperator{\sign}{sgn}
\DeclareMathOperator{\sgn}{sgn}
\DeclareMathOperator{\diag}{diag}
\DeclareMathOperator{\trace}{tr}
\DeclareMathOperator{\arsh}{arsh}

\renewcommand{\phi}{\varphi}
\newcommand{\Studt}{\mathsf{t}}
\newcommand{\Fish}{\mathsf{F}}
\newcommand{\DF}{\ensuremath{\mathsf{DF}}}
\newcommand{\DW}{\ensuremath{\mathsf{DW}}}
\newcommand{\WN}{\ensuremath{\mathsf{WN}}}
\newcommand{\hypo}{\mathcal{H}}
\newcommand{\ARMA}[2]{\ensuremath{\text{ARMA}(#1; #2)}}
\newcommand{\ARCH}[1]{\ensuremath{\text{ARCH}(#1)}}
\newcommand{\GARCH}[2]{\ensuremath{\text{GARCH}(#1; #2)}}
\newcommand{\ARIMA}[3]{\ensuremath{\text{ARIMA}(#1; #2 ; #3)}}
\newcommand{\AR}[1]{\ensuremath{\text{AR}(#1)}}
\newcommand{\MA}[1]{\ensuremath{\text{MA}(#1)}}
\newcommand{\VAR}[1]{\ensuremath{\text{VAR}(#1)}}
\newcommand{\pval}{\ensuremath{p_{\text{val}}}}

\newcommand{\ENGs}[1]{\foreignlanguage{english}{#1}}

\newcommand{\primred}[1]{{\color{red} \textit{Прим. ред.}\quad #1}}

\newcommand{\proglang}[1]{\textsf{#1}}
\newcommand{\pkg}[1]{\textbf{#1}}
\newcommand{\code}{\verb}

\DeclareUnicodeCharacter{0406}{{\fontencoding{X2}\selectfont\CYRII}} % И ДЕС.
\DeclareUnicodeCharacter{0456}{{\fontencoding{X2}\selectfont\cyrii}} % и дес.
\DeclareUnicodeCharacter{0462}{{\fontencoding{X2}\selectfont\CYRYAT}} % ЯТЬ
\DeclareUnicodeCharacter{0463}{{\fontencoding{X2}\selectfont\cyryat}} % ять
\DeclareUnicodeCharacter{0472}{{\fontencoding{X2}\selectfont\CYROTLD}} % ФИТА
\DeclareUnicodeCharacter{0473}{{\fontencoding{X2}\selectfont\cyrotld}} % фита
\DeclareUnicodeCharacter{0474}{{\fontencoding{X2}\selectfont\CYRIZH}} % ИЖИЦА
\DeclareUnicodeCharacter{0475}{{\fontencoding{X2}\selectfont\cyrizh}} % ижица

% % % % % Собственный стиль

\raggedbottom
\righthyphenmin=2
\let\myfootnote\footnote
\renewcommand{\footnote}[1]{\myfootnote{\;#1}}
\setlength{\parindent}{1.5em}
\numberwithin{equation}{section}

\renewcommand{\thesubfigure}{\thefigure.\arabic{subfigure}}
\makeatletter
\renewcommand{\p@subfigure}{}
\renewcommand{\@thesubfigure}{\textbf{Рис. \thesubfigure.\hskip\subfiglabelskip}}
\makeatother

\hyphenation{ап-про-к-си-ми-ро-вать}

\endinput
%%
%% End of file `epsilonj.sty'.


\RequirePackage{graphicx}
\RequirePackage[colorlinks,citecolor=blue,urlcolor=blue]{hyperref}

\newcommand{\specialcell}[2][c]{\begin{tabular}[#1]{@{}c@{}}#2\end{tabular}}

\begin{document}

\TITLE{Кто выигрывает рэп-баттлы?}
\SHORTTI{Кто выигрывает рэп-баттлы?}

\AUTHOR{Андрей Зубанов}{НИУ ВШЭ, Москва.}
\SHORTAU{Зубанов А.}

\thispagestyle{firststyle}

\maketitle

\begin{abstract}
Статистика является мощным инструментом при изучении не только экономических и~социальных, но и~культурных явлений. Даже простые статистические методы могут оказаться полезными при изучении таких событий, как различные соревнования, телешоу, прослушивания.
\end{abstract}

\begin{keyword}
рэп, баттл, жеребьёвка
\end{keyword}



\section{Влияние позиции участника на результат}

Данная статья анализирует рэп-баттлы на примере популярного российского «Versus Battle». В~исследование вошли данные, собранные по выпускам «Первый сезон», «Второй сезон» и~«Межсезонье». В~баттле участвуют два человека (реже "--- две команды). Сущность рэп"=баттла «Versus» состоит в~том, чтобы посредством заранее подготовленного речитатива высказаться о~себе и~своём сопернике. Баттл проходит без музыкального сопровождения. Победитель определяется тремя судьями простым большинством голосов. На протяжении рассматриваемых баттлов участники стоят лицом друг к~другу, сбоку от них стоят судьи, вокруг располагаются немногочисленные болельщики, приглашённые на мероприятие. Баттл состоит из трёх раундов, каждый из которых начинает участник, выбранный жеребьёвкой первым.

\begin{figure}[htbp]
	\centering
	\includegraphics[width=8cm]{maxresdefault.jpg}
	\caption{Рэп-баттл}
\end{figure}

\begin{table}[htb] \centering
		\begin{tabular}{|c|c|c|}
			\toprule
			\specialcell{Победитель \\ по номеру} & \specialcell{Положение  первого \\ относительно судей} & \specialcell{Выиграл ли \\ участник справа} \\
			\midrule
			1& справа&1\\
			1& слева&0\\
			2& слева&1\\
			2& слева&1\\
			1& слева&0\\
			1& слева&0\\
			2& слева&1\\
			2& слева&1\\
			2& слева&1\\
			2& слева&1\\
			2& справа&0\\
			2& справа&0\\
			2& слева&1\\
			2& слева&1\\
			2& слева&1\\
			2& слева&1\\
			2& слева&1\\
			2& справа&0\\
			2& справа&0\\
			1& справа&1\\
			2& слева&1\\
			2& слева&1\\
			1& слева&0\\
			2& слева&1\\
			2& справа&0\\
			2& слева&1\\
			2& справа&0\\
			2& слева&1\\
\bottomrule			
		\end{tabular}
		\caption{Данные о~победах первого/второго по очереди участника и~участника, стоящего слева/справа от судей}
	\end{table}

Согласно данным, всего победил первый по очереди игрок 6~раз, второй игрок "--- 22~раза. Всего победил игрок слева 18~раз, справа "--- 10~раз. 

Основной гипотезой исследования является гипотеза о~том, что второй участник вследствие очерёдности хода имеет преимущество, предположительно из-за возможности ответить в~своём выступлении на выступление соперника. Также это возможно и~в~силу того, что выступающий последним лучше запоминается судьям и~поэтому выше оценивается. Вторая гипотеза "--- предположение о~том, что из-за своего положения относительно судей участник слева имеет преимущество.

Если результат баттла не зависит от очерёдности, то при любом назначенном номере участник должен выигрывать с~вероятностью, близкой к~$1/2$. Проверим эту гипотезу, учитывая, что номер выигравшего участника "--- случайная величина, распределённая биномиально.
\begin{equation}
	\frac{|\hat p_n-p_0|\sqrt{n}}{\sqrt{\hat p_n(1-\hat p_n)}}\sim \Studt_{28}
\end{equation}

Доля выигравших вторых игроков составляет $22/28$, всего наблюдений в~выборке~28. Значит,
\begin{equation}
z_{\text{расч}}=\frac{|22/28-1/2|\sqrt{28}}{\sqrt{22/28 \cdot (1-22/28)}}=8{,}98 > 2{,}05=\Studt_{0{,}975; 28}
\end{equation}

Таким образом, выступающие вторыми участники значимо чаще выигрывают рэп"=баттлы.

Также и~участники слева чаще выигрывают баттлы (вероятность $18/28$).
\begin{equation}
z_{\text{расч}}=\frac{|18/28-1/2|\sqrt{28}}{\sqrt{18/28 \cdot (1-18/28)}}=3{,}29>2{,}05=\Studt_{0{,}975; 28} 
\end{equation}
Однако при уровне значимости в~$0{,}001$ критическая статистика составляет $\Studt_{0{,}0005; 28}=3{,}67$, и~значимого различия не выявляется.

\primred{Кто берёт такие уровни значимости? Зачем?}

\section{Выводы}

Вопрос о~зависимости между позицией участника и~очерёдностью остаётся открытым, однако первое не должно влиять на второе, так как расстановка происходит после жеребьёвки.

Несмотря на то что жеребьёвка не показывается зрителю, а~организаторы не раскрывают процедуры, допустимо считать её случайной, так как это соответствует традициям рэп"=баттлов.

В~итоге нельзя с~уверенностью сказать о~причинно"=следственных связях между очерёдностью участников в~баттлах и~их результатами, ведь может существовать третий фактор, влияющий и~на то, и~на другое. Однако можно утверждать, что вторые участники действительно значимо чаще выигрывают рэп"=баттлы.

\vfill

\begin{center}
		\includegraphics[width=2cm]{versus.jpg}
\end{center}


\end{document}
